% Generated draft: requires mathematical and source review.
\documentclass[11pt]{article}
\usepackage[margin=25mm]{geometry}
\usepackage{fontspec}
\setmainfont{Latin Modern Roman}
\usepackage{amsmath,amssymb,mathtools,unicode-math}
\setmathfont{Latin Modern Math}
\usepackage{enumitem,needspace,xurl,hyperref}
\setlength{\parindent}{0pt}
\setlength{\parskip}{5pt}
\setlength{\emergencystretch}{3em}
\newcommand{\N}{\mathbb{N}}
\newcommand{\Z}{\mathbb{Z}}
\newcommand{\Q}{\mathbb{Q}}
\newcommand{\R}{\mathbb{R}}
\newcommand{\C}{\mathbb{C}}
\newcommand{\F}{\mathbb{F}}
\newcommand{\A}{\mathbb{A}}
\newcommand{\PP}{\mathbb{P}}
\newcommand{\E}{\mathbb{E}}
\newcommand{\eps}{\varepsilon}
\newcommand{\dd}{\,\mathrm{d}}
\begin{document}
\section*{UnsolvedMath ID: 30003803}
\textbf{Source problem code:} catalogue entry
\subsection{Definitions and mathematical statement}
Let \(k\) be a field and let \(S\) be a functor from field extensions of \(k\) to pointed sets. For each extension \(l/k\), denote the distinguished point of \(S(l)\) by \(*_l\). For \(s\in S(k)\), write \(s_l\) for its image under the restriction map \(S(k)\to S(l)\); thus \(s_l\) is trivial precisely when \(s_l=*_l\).

A \(k\)-variety is understood here in the usual sense of a separated scheme of finite type over \(k\). It is geometrically integral if its base change to an algebraic closure of \(k\) is integral. For a \(k\)-variety \(X\) and an extension \(l/k\), put \(X_l=X\times_{\operatorname{Spec}k}\operatorname{Spec}l\) and \(X(l)=\operatorname{Hom}_k(\operatorname{Spec}l,X)\).

\emph{Question as supplied.} For every field \(k\), every functor \(S\) as above, and every \(s\in S(k)\), does there exist a smooth geometrically integral \(k\)-variety \(X\) such that, simultaneously for every field extension \(l/k\),
\[
 s_l=*_l\quad\Longleftrightarrow\quad X(l)\ne\varnothing,
\]
and such that \(X(l)\) is Zariski-dense whenever it is nonempty? The variety \(X\) may depend on \((k,S,s)\), but must be chosen independently of \(l\).

\emph{Unresolved qualifications.} The source says that \(X(l)\) is dense ``in \(X\).'' Taken literally, this means that the images in the underlying topological space of \(X\) of all morphisms in \(X(l)\) have closure \(X\). Another possible intended condition is that the corresponding \(l\)-rational points are Zariski-dense in \(X_l\). The supplied record does not distinguish these formulations, so the ambient space for the density requirement requires confirmation. No additional hypotheses on \(S\), no restriction to infinite fields, and no restriction on the extensions \(l/k\) are supplied. In particular, the quantifiers include finite fields and the extension \(k/k\). Whether the unrestricted functor formulation captures the intended workshop problem requires source review.
\subsection{Short English statement}
For a field \(k\), a pointed functor \(S\), and an element \(s\in S(k)\), the question asks for one smooth geometrically integral variety whose rational points over every extension detect exactly when \(s\) becomes trivial. Whenever such points exist, they must also be Zariski-dense. The source leaves the ambient space for density and any intended restrictions on the functor unresolved. This is a statement-only draft requiring human mathematical review.
\subsection{Sources}
\begin{itemize}
\item Supplied local UnsolvedMath registry record, permanent ID \(30003803\), title \emph{Generic Splitting Varieties for Pointed Functors}, collection \emph{Oberwolfach Reports Open Problems}, exact source problem code \texttt{catalogue entry}. No individual proposer is specified; the supplied proposed year is 2018.
\item Supplied report citation: \emph{Quadratic Forms and Related Structures over Fields (2019)}, Oberwolfach Reports, pp.\ 1235--1291, DOI \texttt{10.4171/owr/2018/21}. \url{https://doi.org/10.4171/owr/2018/21}. The background attributes the problem to workshop open-problem discussions documented in Oberwolfach Reports. The report title's year 2019 and the supplied year 2018 are retained without reconciliation.
\item Supplied literature reference labelled ``Demarche and Florence on norm varieties'': \url{https://arxiv.org/abs/1711.06585}.
\item Supplied literature reference labelled ``Rost norm varieties'': \url{https://doi.org/10.1016/j.jpaa.2005.12.012}.
\item Literature assessment dated exactly \texttt{2026-08-21}: the supplied curation labels the problem open and states that Rost norm varieties give the required criterion only up to a zero cycle of degree prime to the symbol prime, while the general one-variety rational-point criterion remains unresolved. This is supplied dated curation derived from the OWR record, not an independently verified literature assessment. The record expressly requests independent verification of time-sensitive status and cited claims.
\item The supplied statement status is \texttt{corrected\_verified}; its accompanying note reports manual semantic, typographic, and LaTeX review. These are registry assertions and do not constitute independent mathematical verification of this draft.
\end{itemize}
\end{document}
